% Bien dich duoc bang pdfLaTeX tren Overleaf.
\documentclass[aspectratio=169,10pt]{beamer}

\usetheme{Madrid}
\setbeamertemplate{navigation symbols}{}

\usepackage[utf8]{inputenc}
\usepackage[T5]{fontenc}
\usepackage[vietnamese]{babel}
\usepackage{lmodern}
\usepackage{booktabs}
\usepackage{tabularx}
\usepackage{array}

\newcommand{\code}[1]{\texttt{\detokenize{#1}}}
\newcolumntype{Y}{>{\raggedright\arraybackslash}X}

\title[Tiny Shell]{Tiny Shell trên Windows}
\subtitle{Project môn Nguyên lý hệ điều hành}
\author{Nhóm thực hiện}
\institute{IT3070}
\date{Tháng 6 năm 2026}

\begin{document}

\begin{frame}{Giới thiệu project}
  \begin{block}{Mục tiêu}
    Xây dựng một shell tối giản trên Windows để minh họa các thao tác cơ bản của hệ điều hành đối với tiến trình.
  \end{block}

  \begin{itemize}
    \item Tên chương trình: \code{myShell}.
    \item Ngôn ngữ: C++17; hệ build: CMake.
    \item Nền tảng: Windows 10/11.
    \item API chính: Windows Process API, ToolHelp API, Console Control Handler và Environment API.
    \item Phạm vi: đọc lệnh, phân tích lệnh, tạo tiến trình con và quản lý tiến trình nền.
  \end{itemize}
\end{frame}

\begin{frame}{Nhóm thực hiện}
  \begin{table}
    \large
    \begin{tabular}{@{}ll@{}}
      \toprule
      Họ và tên & MSSV \\
      \midrule
      Nguyễn Nhật Minh & 202416291 \\
      Bùi Quang Phúc & 202400067 \\
      Nguyễn Công Hiếu {\color{blue}(C)} & 202416202 \\
      Phùng Danh Chí Vĩ & 202416398 \\
      \bottomrule
    \end{tabular}
  \end{table}
\end{frame}

\begin{frame}{Tính năng chính}
  \begin{table}
    \small
    \begin{tabularx}{\textwidth}{@{}lY@{}}
      \toprule
      Nhóm chức năng & Nội dung kỹ thuật \\
      \midrule
      Command parsing & Phân tích command line thành tokens, xử lý dấu nháy kép và nhận diện background operator \code{&}. \\
      Process creation & Tạo child process từ external command bằng Windows API \code{CreateProcessA()}. \\
      Foreground execution & Với foreground process, shell gọi \code{WaitForSingleObject()} và ghi nhận exit code sau khi process kết thúc. \\
      Background execution & Với command có \code{&}, shell không wait mà lưu process vào job table và trả prompt cho user. \\
      Job control & Theo dõi PID, \code{PROCESS_INFORMATION}, status và exit code; hỗ trợ \code{list}, \code{kill}, \code{stop}, \code{resume}, \code{reap}. \\
      Script và batch support & Chạy script \code{.tsh}; chuyển \code{.bat}/\code{.cmd} qua \code{cmd.exe /C} theo cơ chế của Windows. \\
      \bottomrule
    \end{tabularx}
  \end{table}
\end{frame}

\begin{frame}{Các lệnh chính}
  \begin{table}
    \small
    \begin{tabularx}{\textwidth}{@{}lY@{}}
      \toprule
      Lệnh & Chức năng \\
      \midrule
      \code{help} & Hiển thị danh sách lệnh được hỗ trợ. \\
      \code{exit} & Thoát shell và dọn các tiến trình nền còn được theo dõi. \\
      \code{list} hoặc \code{jobs} & Liệt kê PID, trạng thái và command của job nền. \\
      \code{kill <pid>} & Kết thúc một tiến trình nền đang được shell quản lý. \\
      \code{stop <pid>} & Tạm dừng các thread của tiến trình nền. \\
      \code{resume <pid>} & Tiếp tục các thread đã bị tạm dừng. \\
      \code{reap} & Xóa các job đã kết thúc khỏi bảng theo dõi. \\
      \code{date}, \code{time}, \code{dir} & Cung cấp một số tiện ích hệ thống cơ bản. \\
      \bottomrule
    \end{tabularx}
  \end{table}
\end{frame}

\begin{frame}{Cơ chế thực thi}
  \begin{itemize}
    \item Shell hoạt động theo chu trình: đọc lệnh, phân tích cú pháp, phân loại lệnh, sau đó thực thi.
    \item Lệnh built-in được xử lý trực tiếp trong tiến trình \code{myShell}.
    \item Lệnh bên ngoài được tạo thành tiến trình con bằng \code{CreateProcessA()}.
    \item Với foreground process, shell lưu tiến trình hiện tại và chờ đến khi tiến trình kết thúc.
    \item Với background process, shell lưu thông tin vào bảng job và tiếp tục nhận lệnh mới.
    \item Ctrl+C được xử lý bằng \code{SetConsoleCtrlHandler()} để hủy foreground process nhưng không đóng shell.
  \end{itemize}
\end{frame}

\begin{frame}{Chạy demo}
\end{frame}

{
\usebackgroundtemplate{\color{blue!4}\rule{\paperwidth}{\paperheight}}
\begin{frame}[plain]
  \centering
  \vfill

  {\color{blue!65!black}\Huge Cảm ơn thầy và các bạn}

  \vspace{3mm}
  {\color{blue!65!black}\Huge đã lắng nghe}

  \vspace{8mm}
  {\color{blue!45!black}\rule{0.56\textwidth}{0.5pt}}

  \vspace{8mm}
  {\LARGE Câu hỏi và thảo luận}

  \vfill
  {\small Tiny Shell trên Windows -- IT3070}
\end{frame}
}

\end{document}
